\chapter{Experiments}
\label{ch:experiments}

This chapter states what is being asked, on what data, against what, and how the
answers will be judged. It is written before the results so that every quantity
in Chapter~\ref{ch:results} has a declared meaning and a declared success
criterion rather than one chosen after the numbers arrived.

\section{Research Questions}
\label{sec:exp-rq}

Four questions. Each is paired with the outcome that would \emph{refute} it,
because a question with no refuting outcome is a name for a group of
experiments rather than a question.

\begin{description}[leftmargin=0pt,style=unboxed,itemsep=10pt]

\item[\textbf{RQ1 --- Does cost-aware validation scheduling reduce the cost of
rejecting a patch?}]
Does ranking tests by $K/\hat{m}^{\,c}$ --- an online rejection rate over an
online cost --- reduce the CPU spent on validation, compared with the order the
benchmark ships, an uninformed random order, the published repair-loop
prioritiser \frtp~\citep{qi2013frtp}, and the classical additional-greedy
coverage prioritiser~\citep{rothermel2001prioritizing}?
\par\smallskip
\emph{Refuted when} the scheduler is not cheaper than the shipped order beyond
the bootstrap interval; or when it is cheaper than the shipped order but not
cheaper than \frtp\ and coverage, in which case the contribution is early
stopping, which is not ours.

\item[\textbf{RQ2 --- Do selected hints reduce redundant generation?}]
Does handing the model a small, freshly selected set of concrete requirements
move its proposals off the ideas it has already tried --- fewer duplicate
patches, and at least as many repairs at the same budget --- where independent
repeated sampling, a prohibition, and a full conversation transcript do not?
\par\smallskip
\emph{Refuted when} the duplicate rate does not fall against independent
repeated sampling; or when it falls and the repair rate does not rise, in which
case the distribution moved somewhere no better.

\item[\textbf{RQ3 --- Do the two mechanisms reinforce each other?}]
The two mechanisms are claimed to compose: scheduling produces the failure
record that hints are selected from, and better hints send fewer wasted
candidates to the expensive validation step. Does the combined system beat the
sum of what each does alone, on total cost and on repairs within budget?
\par\smallskip
\emph{Refuted when} the combination is no better than the better of the two
alone, in which case the mechanisms are independent and should be reported as
two separate contributions rather than one system.

\item[\textbf{RQ4 --- Which parts carry the effect, and how sensitive are they?}]
Three things separate the proposed scheduler from the shipped order: having a
memory at all, turning a count of rejections into a rate, and dividing that rate
by cost. Which one produces the effect? And does the answer survive the design
choices --- the cost exponent $c$, the cold-start rule, the exploration rate,
and the number of hints shown?
\par\smallskip
\emph{Refuted when} no single step clears the bootstrap interval, so the gain is
simply ``having any memory'' and the score function is ornament; or when the
answer flips under a routine change of one parameter, in which case the result
is an artefact of tuning.

\end{description}

\section{Benchmark: ConDefects Python}
\label{sec:exp-benchmark}

\subsection{Corpus, and the 100-Task Pilot Drawn From It}

The corpus is the Python subset of ConDefects~\citep{wu2024condefects}, chosen
for the four reasons given in Section~\ref{sec:bg-benchmark}: every test input ships
with its expected output, a test is one whole-program execution, the inputs
within a single task are deliberately heterogeneous in cost, and the faults
post-date the training cut-offs of commonly benchmarked models.

The evaluation reported here is a \textbf{pilot of $100$ tasks} drawn from that
corpus; the scale and its reasons are in
Section~\ref{sec:exp-pilot-scale} below, and running the full corpus is the
first item of future work (Section~\ref{sec:future}).

A task is admitted to the corpus when every one of the following holds.
The screen is applied before any configuration runs, and the admitted set is
frozen with a checksum so that no later analysis can change it.

\begin{table}[htbp]
\centering
\caption{Corpus admission criteria. Applied once, before treatment.}
\label{tab:gate}
\small
\begin{tabularx}{\textwidth}{@{}l X@{}}
\toprule
Criterion & Why\\
\midrule
Single-file Python fault & The scope of this study; matches the sandbox model
of one program, one input.\\
Faulty and reference versions both present & The reference supplies the expected
outputs the oracle compares against.\\
At least one test rejects the faulty version & Otherwise the task is already
``repaired'' at round~$0$ and measures nothing.\\
The reference passes every test in the pool & A test the correct program fails
is an instrument fault, not a fault in the candidate; such tests leave the
pool.\\
The reference finishes within the sandbox clock on every pooled test & Applying
the clock to the reference first is what keeps the oracle sound: a test the
correct program cannot finish is not a test at all.\\
Pool size at least the declared minimum & A scheduling experiment on a pool of
three tests measures nothing about scheduling.\\
\bottomrule
\end{tabularx}
\end{table}

\pendnote{The admission funnel --- how many faults each criterion removed, and
how many tasks the $100$ were drawn from --- is reported in
Section~\ref{sec:res-corpus}.}

\subsection{Pilot Scale and Why It Is a Pilot}
\label{sec:exp-pilot-scale}

\begin{table}[htbp]
\centering
\caption{Pilot configuration.}
\label{tab:pilot-config}
\small
\begin{tabularx}{\textwidth}{@{}l l X@{}}
\toprule
Parameter & Value & Note\\
\midrule
Tasks            & $100$            & Drawn from the admitted ConDefects Python corpus\\
Seeds per task   & $3$              & Different draws from the same model, same task\\
Budget $B$       & $20$ rounds      & Maximum candidate patches per session\\
Runs per configuration & $6{,}000$ rounds & $100 \times 3 \times 20$ at the budget ceiling\\
Sandbox clock    & $10$\,s wall     & A hard kill, applied to the reference first\\
Proposer model   & \code{qwen2.5-\hspace{0pt}coder:\hspace{0pt}7b} & One open-weights model served locally
through an OpenAI-compatible endpoint, held fixed across every configuration\\
Temperature      & $1.0$            & Fixed; the study varies policy, not decoding\\
\bottomrule
\end{tabularx}
\end{table}

This is a pilot and is reported as one. Its purpose is to establish whether the
effects exist and are large enough to be worth a corpus-scale run, and to fix
the parameters that the full run must declare in advance. The full ConDefects
Python corpus is roughly an order of magnitude larger, and running it is the
next step (Section~\ref{sec:future}).

\subsection{Screens Are Bands, Decided After the Run}
\label{sec:exp-bands}

One protocol decision governs the whole corpus and is worth stating before the
configurations, because it determines what can be asked afterwards.

A \emph{screen} rejects a task before any measurement exists; a \emph{band}
keeps every task and groups them at analysis time. They are not equally
reversible --- a screen must be decided before the first cell exists and cannot
be undone, whereas a band can be re-cut any number of ways from the same rows
--- so where both are available this study takes the reversible one. The
corpus admits every task whose oracle is sound: the criteria in
Table~\ref{tab:gate} are only those without which a measured number would be
\emph{wrong}, never those that merely shape the population.

Table~\ref{tab:bands} declares the three band axes before the corpus run, one
per condition of Section~\ref{sec:disc-limits}.

\begin{table}[htbp]
\centering
\caption[Band axes, declared before the run]{The three band axes, their
boundaries, and when each becomes known. Boundaries are fixed here so that the
banded result in Chapter~\ref{ch:results} is not a cut chosen after seeing
which one looks best.}
\label{tab:bands}
\small
\begin{tabularx}{\textwidth}{@{}>{\raggedright\arraybackslash}p{0.14\textwidth} >{\raggedright\arraybackslash}p{0.34\textwidth} X@{}}
\toprule
Axis & Bands & Status\\
\midrule
Pool size $|P|$ &
\code{degenerate} $=1$; \code{tiny} $2$--$9$; \code{small} $10$--$19$;
\code{standard} $\ge 20$ &
A property of the frozen pool, known before any candidate is drawn. A
\emph{stratification}\\
\addlinespace
Within-task cost spread &
\code{flat} / \code{mixed} / \code{spread}, by terciles of the within-task
coefficient of variation of per-test CPU &
A property of the frozen cost table, established by the arm-independent
\code{audit} pass before any configuration runs. A \emph{stratification}\\
\addlinespace
Task difficulty &
\code{dead} $=0$; \code{hard} $<0.25$; \code{mid} $0.25$--$0.60$;
\code{easy} $0.60$--$1$; \code{trivial} $=1$ &
The repair rate of the floor configuration, per model. An \textbf{outcome}, so
it may only ever be a post-hoc band and never an admission criterion\\
\bottomrule
\end{tabularx}
\end{table}

Two rules keep the difficulty axis honest, because it is the one that could be
abused. The rate is taken from a single designated floor configuration
(\code{no-mem}), not from the configuration under test, so no policy can band
itself into a favourable population. And it is frozen once, beside the corpus
checksum, and never recomputed by an analysis script --- a rate recomputed on
every call is a rate that moves when the merge moves, and every banded number
downstream would move with it without saying so.

The \code{standard} pool band is exactly what a pre-run screen of $|P| \ge 20$
would have admitted, so every number such a screen would have produced is still
computable, and the tasks it would have deleted are reported beside it rather
than discarded.

\section{Baselines and Configurations}
\label{sec:exp-config}

\subsection{The Two Blocks}

Configurations fall into two blocks that differ in what they cost to run, and
the difference governs the order of work.

\textbf{Block A} varies only the \emph{order} in which tests are run. The
prompt is identical across these configurations, so the candidate patches can be
drawn once and reused: adding a Block~A configuration costs CPU but no model
calls.

\textbf{Block B} varies the \emph{prompt}. Here the round-$t$ prompt depends on
what validation found in rounds $1..t-1$, so the draw and the schedule cannot be
separated and each configuration needs its own set of model calls.

One identity ties the two blocks together. In round~$1$ the store is empty,
so a steered prompt carries no hints and is byte-identical to an unsteered one.
Round~$1$ must therefore be the same draw in every configuration, which requires
the draw seed to be a function of the task, the seed and the round number only
--- never of the policy. Deriving it from the full configuration key would make
every configuration buy its own round~$1$ and would silently destroy the pairing
that every comparison in Section~\ref{sec:exp-stats} depends on.

\subsection{Block A --- Scheduling Configurations}

\begin{table}[htbp]
\centering
\caption{Block A: configurations that differ only in test order. ``Answers''
names the research question each one serves.}
\label{tab:blocka}
\small
\begin{tabularx}{\textwidth}{@{}>{\raggedright\arraybackslash}p{0.215\textwidth} X l@{}}
\toprule
Name & Order of tests & Answers\\
\midrule
\code{full-suite}    & Every test, no early stop & calibration\\
\code{no-mem} \newline \emph{(independent} \newline \emph{repeated sampling)}
                     & Shipped order, and the store is never updated, so every round
                       redraws from scratch. The name is the implementation's; what it
                       \emph{is} is independent repeated prompting & calibration, RQ2 baseline\\
\code{natural}       & The order the benchmark ships --- \textbf{the incumbent} & RQ1\\
\code{random}        & Seeded shuffle, reshuffled each round --- \textbf{the control} & RQ1\\
\code{frtp}          & Proven rejecters by \emph{count} descending, then the rest~\citep{qi2013frtp} & RQ1, RQ4\\
\code{coverage}      & Additional-greedy over reference line coverage~\citep{rothermel2001prioritizing} & RQ1\\
\code{cost-only}     & $\hat{m}$ ascending; no rejection information & RQ4\\
\code{kill-only}     & $K$ descending; no cost information & RQ4\\
\textbf{\code{sched}} & $K/\hat{m}^{\,c}$ --- \textbf{the proposed scheduler} & RQ1, RQ3\\
\midrule
\code{sched-cs-median} \newline \code{sched-cs-q75} & Cold start at the median or the $75$th percentile instead of the max & sensitivity\\
\code{sched-bucket}  & Discrete buckets instead of a continuous score & sensitivity\\
\code{sched-explore0}& $\varepsilon = 0$ & sensitivity, RQ4 precondition\\
\bottomrule
\end{tabularx}
\end{table}

Two framing points govern how this table is read.

\paragraph{\code{natural} is the incumbent, not a neutral control.} Sorting by
file name is not a scheduling decision anyone made; it is what any
implementation gets for free, and therefore what most implementations do. It is
also not arbitrary: contest naming puts the small, public, highly discriminating
sample cases first, so the shipped order is \emph{better} than chance. The
uninformed control is \code{random}, and reporting both is what separates
``better than nothing'' from ``better than the thing people actually run''.

\paragraph{\code{frtp} is the published baseline, implemented as published.} It
ranks on the rejection \emph{count}, which is what \frtp\ ranks on, not on a
rate and not on anything divided by cost. A baseline weakened to make the
proposal look better answers nothing.

\subsection{Block B --- Prompt Configurations}

\begin{table}[htbp]
\centering
\caption{Block B: configurations that differ in what the prompt carries. All
use the proposed scheduler unless the name says otherwise.}
\label{tab:blockb}
\small
\begin{tabularx}{\textwidth}{@{}l X l@{}}
\toprule
Name & What the prompt carries after round~1 & Answers\\
\midrule
\code{hints}          & $k{=}2$ selected requirements, \textbf{shipped test order} & RQ2, RQ3\\
\code{sched+hints}    & $k{=}2$ selected requirements \textbf{and} the proposed scheduler & RQ2, RQ3\\
\code{exclusion}      & The same facts phrased as a prohibition & RQ2\\
\code{chat}           & The full conversation transcript, \textsc{ChatRepair}-style~\citep{xia2024chatrepair} & RQ2\\
\midrule
\code{hints-k1}, \code{hints-k3} & $k \in \{1,3\}$ instead of $2$ & RQ4\\
\code{hints-sample}   & Disclosure \code{sample-only} & RQ4\\
\code{hints-noreq}    & Disclosure \code{no-required-line} & RQ4\\
\bottomrule
\end{tabularx}
\end{table}

\paragraph{\code{no-mem} is the baseline for RQ2, and it is a strong one.} It
redraws independently every round under a policy-free seed, which is exactly
\emph{independent repeated prompting}. \citet{chen2025conversational} measured
that strategy at $23$ bugs repaired against \textsc{ChatRepair}'s $18$ at an
equal budget of $24$ model calls, with $65\,\%$ of the conversational system's
patches being duplicates. The obvious iterative strategy is therefore already
published as \emph{worse} than simply drawing again. Block~B must beat
\code{no-mem}; beating \code{natural} proves nothing on the generation side.

\subsection{The Factorial for RQ3}

RQ3 asks whether the two mechanisms reinforce each other, which requires each of
them alone and both together on the same tasks and seeds.
Table~\ref{tab:factorial} is that design.

\begin{table}[htbp]
\centering
\caption{The $2\times2$ design behind RQ3. Each cell is a configuration run on
the same $100$ tasks $\times$ $3$ seeds.}
\label{tab:factorial}
\small
\begin{tabular}{@{}l l l@{}}
\toprule
 & \textbf{Shipped test order} & \textbf{Cost-aware scheduling}\\
\midrule
\textbf{No hints}       & \code{no-mem} (baseline) & \code{sched}\\
\textbf{Selected hints} & \code{hints}             & \textbf{\code{sched+hints}} (full system)\\
\bottomrule
\end{tabular}
\end{table}

\subsection{Parameter Sweeps}

Two parameters are measured rather than assumed. The \textbf{cost exponent} $c$
of Equation~\eqref{eq:score} is swept over the range $[0, 1.25]$, whose two
endpoints are the two published positions: $c=0$ is ranking on the rejection
rate alone, $c=1$ is full cost-cognizance. The \textbf{hint dose} is swept over
$k \in \{1,2,3\}$ crossed with the three disclosure levels of
Table~\ref{tab:disclosure}.

Both sweeps are reported as sensitivity, not as headline results. A value
selected on the same data it is reported on is optimistically biased, so the
full corpus run must declare $c$ and $k$ in advance and report the sweep as a
check (Section~\ref{sec:future}).

\section{Metrics}
\label{sec:exp-metrics}

The study does not collapse into a single figure of merit, because the two
halves of the loop spend different resources and a composite would let one pay
for the other.

\begin{table}[htbp]
\centering
\caption{Validation-side metrics (RQ1, RQ3, RQ4).}
\label{tab:metrics-validation}
\small
\begin{tabularx}{\textwidth}{@{}>{\raggedright\arraybackslash}p{0.21\textwidth} X X@{}}
\toprule
Metric & Definition & Success criterion\\
\midrule
CPU seconds per session & \code{getrusage} delta summed over the tests actually
run; paired per task and seed & \textbf{Primary.} The scheduler below
\code{natural}, \code{random}, \code{frtp} and \code{coverage}, beyond the
bootstrap interval\\
Test executions per session & How many tests were run; reported because the
prioritisation literature reports executions & Reported for every configuration;
\emph{not} a success criterion on its own\\
CPU per execution & CPU divided by executions; separates ``cheaper'' from
``fewer'' & Diagnostic: it is where the mechanism claim lives\\
Stepwise decomposition & Paired $\Delta$ CPU for each one-variable step of the
ablation ladder (Table~\ref{tab:ladder}) & At least one step clears the
interval\\
\bottomrule
\end{tabularx}
\end{table}

\begin{table}[htbp]
\centering
\caption{Generation-side metrics (RQ2, RQ3, RQ4).}
\label{tab:metrics-generation}
\small
\begin{tabularx}{\textwidth}{@{}>{\raggedright\arraybackslash}p{0.21\textwidth} X X@{}}
\toprule
Metric & Definition & Success criterion\\
\midrule
Duplicate rate & $1 - \text{distinct}/\text{total}$ patches in a session, over
normalised patch text & \textbf{Primary for RQ2.} Falls against \code{no-mem}\\
\code{success@B} & Share of sessions that reach a plausible patch within budget
$B$ & Rises, or at minimum does not fall, wherever the duplicate rate falls\\
First-accept round & The round at which the first plausible patch appears &
Reported paired\\
Prompt and completion tokens & The generation bill, since the LLM-repair
literature prices repair in tokens & Reported for every Block~B configuration\\
Hint leakage & Of the patches that are rejected at all, the share whose
rejecting tests are \emph{disjoint} from what the prompt showed --- the patch
satisfied everything it was told and died elsewhere. Computed from the
\code{facts\_shown} field on every round record against the full rejection
vector produced by the \code{audit} pass & \textbf{Primary for the RQ4 dose
curve.} Should fall monotonically as disclosure rises, or the disclosure axis
does not exist\\
\bottomrule
\end{tabularx}
\end{table}

\paragraph{Hint leakage needs no held-out tests.} The leakage metric is
computed entirely from data the main grid already produces: each round record
stores the facts the prompt showed, and the \code{audit} full-scan pass stores,
for every patch, the complete set of tests that reject it. The share of
rejections falling outside the disclosed set is then a set difference over
those two, so the measurement works at the default setting where the whole pool
is the oracle and nothing is held out.

A \emph{second}, narrower question does need held-out tests: of the patches the
loop \emph{accepted}, how many fail on tests the loop never ran? With no
holdout, acceptance is defined by the same tests that would have to judge it,
so no post-hoc analysis can recover that number. That question is answered by a
separate pass with \code{--holdout} on a \emph{subset} of tasks; it is a
different question from RQ1 and does not need the full grid, so it does not
require re-freezing the pilot corpus.

\paragraph{Guaranteed, not measured.} By Definition~\ref{def:soundness},
\code{success@B} and the first-accept round are identical across every Block~A
configuration, on every task and seed, not merely on average. The experiment
therefore \emph{asserts} that equality and halts if it is violated, rather than
reporting it as a finding. A difference there would be a soundness bug.

\section{The Ablation Ladder}
\label{sec:exp-ladder}

An ablation chain must change one thing per step or it cannot attribute the
effect. Going directly from \code{frtp} to \code{sched} changes two things at
once --- a count becomes a rate, \emph{and} the rate is divided by cost --- so
the ladder is built to separate them.

\begin{table}[htbp]
\centering
\caption{The ablation ladder for RQ4. One variable per step.}
\label{tab:ladder}
\small
\begin{tabularx}{\textwidth}{@{}l l X@{}}
\toprule
Step & What changes & What it means if this step does not clear the interval\\
\midrule
\code{natural} $\to$ \code{cost-only}  & Cost enters; no rejection information &
Cost alone is not a usable signal\\
\code{natural} $\to$ \code{kill-only}  & Rejection rate enters; no cost &
The memory term is ornament\\
\code{natural} $\to$ \code{frtp}       & The published baseline's step &
\frtp\ does not reproduce in this setting\\
\code{frtp} $\to$ \code{kill-only}     & A \textbf{count} becomes a \textbf{rate} &
Reading the denominator buys nothing\\
\code{kill-only} $\to$ \code{sched}    & The rate is divided by cost &
Cost-awareness is ornament\\
\code{sched} $\to$ \kmc{c}             & The exponent is tuned &
$c$ need not be a parameter\\
\bottomrule
\end{tabularx}
\end{table}

\section{Implementation and Measurement}
\label{sec:exp-impl}

\subsection{The Tool}

\caster\ is a Python package of command-line verbs; the ones needed to
reproduce the pilot are \code{freeze} (apply the admission screen and write the
frozen corpus), \code{propose} (fill the response cache and patch store),
\code{audit} (a full-scan pass establishing, for every patch, which tests
reject it and what each costs), \code{evaluate} (walk the task $\times$ seed
$\times$ configuration grid, resumably), then \code{merge}, \code{analyse} and
\code{report}. Three module boundaries are asserted by unit tests: only the
benchmark adapter reads the benchmark, the store has no file access and no
benchmark imports, and the prompt module does not import the validator. The
package ships $263$ unit tests that run without a network, a model or the
benchmark.

\subsection{How Cost Is Measured}

Cost is the delta of \code{getrusage(RUSAGE\_CHILDREN)} across one test run:
CPU seconds, never wall clock. Wall clock on a machine running several workers
is contaminated by siblings, which would make the learned cost estimate a
function of machine load rather than of the test. Because \code{getrusage}
accounting is per process, parallelism is process-per-worker and the sandbox
refuses to run under a thread pool rather than let one worker's CPU be
attributed to another.

Three further rules keep the measurement honest.

\begin{enumerate}[leftmargin=*,itemsep=4pt]
\item \textbf{The cost table is established by the audit pass, not by any
configuration.} If a configuration were allowed to populate the cost table, that
table would become a function of which configuration ran first, and every later
comparison would inherit the bias.
\item \textbf{No (patch, test) pair is executed twice.} Outcomes are memoised,
and the identity
$(\text{fresh runs} - \text{timeout re-runs}) + \text{cache hits} = \sum
\text{test calls}$
is instrumented per scan and over the whole run rather than asserted in prose. A
stored outcome that disagrees with a fresh one halts the run.
\item \textbf{Timeouts are quarantined.} A timed-out test is never served from
the cache and never enters the cost model, because a correct-but-slow program
times out too and entering timeouts at the clock value would let a handful of
them dominate the cost estimate.
\end{enumerate}

\subsection{Reproducibility}

The admitted corpus, the pool order and the parameter values are frozen before
any configuration runs, under a checksum. Every value that can change a
measured quantity is either part of a cache key --- model, temperature, prompt,
output budget, draw seed --- or recorded in the run's protocol record, which
carries the sandbox clock, the interpreter version, the backend and the model
digest. Merging results across machines refuses to join across a disagreement
in that record, so a parameter changed halfway through a run produces a refusal
rather than a quietly mixed average.

\section{Statistical Analysis}
\label{sec:exp-stats}

\paragraph{Pairing.} Every comparison is paired at the level of a
(task, seed) cell under common random numbers, so that two configurations see
the same round-$1$ candidate on the same task. Unpaired comparison on $100$
tasks with the cost spread this corpus has would be dominated by which tasks
landed in which arm.

\paragraph{Intervals.} Per-task totals are not independent across seeds, so
confidence intervals come from a \emph{cluster bootstrap over tasks} ---
resampling whole tasks with replacement, $10{,}000$ resamples --- rather than
from a bootstrap over cells, which would treat three seeds on one task as three
independent observations and report an interval that is too narrow.

\paragraph{Tests.} Paired differences are tested with the Wilcoxon signed-rank
test over per-task totals. The distribution of per-task CPU is heavily
right-skewed, so a $t$-test on the mean would be driven by a few expensive
tasks.

\paragraph{Effect size.} Differences are reported as a percentage change in
total CPU together with the number of tasks on which the comparison goes each
way. The second number matters: an aggregate win driven by three expensive tasks
is a different claim from a win on eighty tasks, and the per-task scatter in
Chapter~\ref{ch:results} shows which it is.

\paragraph{Multiple comparisons.} Four primary contrasts are declared in
advance: \code{sched} against \code{natural}, \code{sched} against \code{frtp}
(RQ1), \code{sched+hints} against \code{no-mem} on duplicate rate (RQ2), and
\code{sched+hints} against the better of \code{sched} and \code{hints} (RQ3).
These four carry a Holm correction. Every other contrast in
Chapter~\ref{ch:results} --- the ablation ladder, both sweeps, and the
sensitivity table --- is reported with its interval and labelled exploratory,
and no claim rests on one of them alone.
